
\documentclass[openany, amssymb, psamsfonts]{amsart}
\usepackage{mathrsfs,comment}
\usepackage[usenames,dvipsnames]{color}
\usepackage[normalem]{ulem}
\usepackage{url}
\usepackage[all,arc,2cell]{xy}
\UseAllTwocells
\usepackage{enumerate}
%%% hyperref stuff is taken from AGT style file
\usepackage{hyperref}  
\hypersetup{%
  bookmarksnumbered=true,%
  bookmarks=true,%
  colorlinks=true,%
  linkcolor=blue,%
  citecolor=blue,%
  filecolor=blue,%
  menucolor=blue,%
  pagecolor=blue,%
  urlcolor=blue,%
  pdfnewwindow=true,%
  pdfstartview=FitBH}   
  
\let\fullref\autoref
%
%  \autoref is very crude.  It uses counters to distinguish environments
%  so that if say {lemma} uses the {theorem} counter, then autrorefs
%  which should come out Lemma X.Y in fact come out Theorem X.Y.  To
%  correct this give each its own counter eg:
%                 \newtheorem{theorem}{Theorem}[section]
%                 \newtheorem{lemma}{Lemma}[section]
%  and then equate the counters by commands like:
%                 \makeatletter
%                   \let\c@lemma\c@theorem
%                  \makeatother
%
%  To work correctly the environment name must have a corrresponding 
%  \XXXautorefname defined.  The following command does the job:
%
\def\makeautorefname#1#2{\expandafter\def\csname#1autorefname\endcsname{#2}}
%
%  Some standard autorefnames.  If the environment name for an autoref 
%  you need is not listed below, add a similar line to your TeX file:
%  
%\makeautorefname{equation}{Equation}%
\def\equationautorefname~#1\null{(#1)\null}
\makeautorefname{footnote}{footnote}%
\makeautorefname{item}{item}%
\makeautorefname{figure}{Figure}%
\makeautorefname{table}{Table}%
\makeautorefname{part}{Part}%
\makeautorefname{appendix}{Appendix}%
\makeautorefname{chapter}{Chapter}%
\makeautorefname{section}{Section}%
\makeautorefname{subsection}{Section}%
\makeautorefname{subsubsection}{Section}%
\makeautorefname{theorem}{Theorem}%
\makeautorefname{thm}{Theorem}%
\makeautorefname{cor}{Corollary}%
\makeautorefname{lem}{Lemma}%
\makeautorefname{prop}{Proposition}%
\makeautorefname{pro}{Property}
\makeautorefname{conj}{Conjecture}%
\makeautorefname{defn}{Definition}%
\makeautorefname{notn}{Notation}
\makeautorefname{notns}{Notations}
\makeautorefname{rem}{Remark}%
\makeautorefname{quest}{Question}%
\makeautorefname{exmp}{Example}%
\makeautorefname{ax}{Axiom}%
\makeautorefname{claim}{Claim}%
\makeautorefname{ass}{Assumption}%
\makeautorefname{asss}{Assumptions}%
\makeautorefname{con}{Construction}%
\makeautorefname{prob}{Problem}%
\makeautorefname{warn}{Warning}%
\makeautorefname{obs}{Observation}%
\makeautorefname{conv}{Convention}%


%
%                  *** End of hyperref stuff ***

%theoremstyle{plain} --- default
\newtheorem{thm}{Theorem}[section]
\newtheorem{cor}{Corollary}[section]
\newtheorem{prop}{Proposition}[section]
\newtheorem{lem}{Lemma}[section]
\newtheorem{prob}{Problem}[section]
\newtheorem{conj}{Conjecture}[section]
%\newtheorem{ass}{Assumption}[section]
%\newtheorem{asses}{Assumptions}[section]

\theoremstyle{definition}
\newtheorem{defn}{Definition}[section]
\newtheorem{ass}{Assumption}[section]
\newtheorem{asss}{Assumptions}[section]
\newtheorem{ax}{Axiom}[section]
\newtheorem{con}{Construction}[section]
\newtheorem{exmp}{Example}[section]
\newtheorem{notn}{Notation}[section]
\newtheorem{notns}{Notations}[section]
\newtheorem{pro}{Property}[section]
\newtheorem{quest}{Question}[section]
\newtheorem{rem}{Remark}[section]
\newtheorem{warn}{Warning}[section]
\newtheorem{sch}{Scholium}[section]
\newtheorem{obs}{Observation}[section]
\newtheorem{conv}{Convention}[section]

%%%% hack to get fullref working correctly
\makeatletter
\let\c@obs=\c@thm
\let\c@cor=\c@thm
\let\c@prop=\c@thm
\let\c@lem=\c@thm
\let\c@prob=\c@thm
\let\c@con=\c@thm
\let\c@conj=\c@thm
\let\c@defn=\c@thm
\let\c@notn=\c@thm
\let\c@notns=\c@thm
\let\c@exmp=\c@thm
\let\c@ax=\c@thm
\let\c@pro=\c@thm
\let\c@ass=\c@thm
\let\c@warn=\c@thm
\let\c@rem=\c@thm
\let\c@sch=\c@thm
\let\c@equation\c@thm
\numberwithin{equation}{section}
\makeatother

%Self notation
%%Notations for Project
\newcommand{\Cat}{\mathbf{Cat}} %the 2-category of small category
\newcommand{\SymMon}{\mathbf{SymMon}} %2-category of sym mon cat, with lax functors and natural transforms
\newcommand{\PsAlg}{\mathbf{PsAlg}} %the pseudo-alg 
\newcommand{\PsFun}{\mathbf{PsFun}} %the pseudo-functor
\newcommand{\vs}{\mathrm{vs}} %very special
\newcommand{\Sets}{\mathbf{Set}} %cat of sets
\newcommand{\Q}{\mathscr{Q}} %Q operad (used for Q_G)
\newcommand{\F}{\mathscr{F}} %F cat of finite sets (used for F_G)
\newcommand{\cP}{\mathscr{P}} %P operad (used for P_G also), stands for curly P
\newcommand{\cO}{\mathscr{O}} %operad in general
%used for $G$-categories
\newcommand{\M}{\mathscr{M}} %associativity operad
\newcommand{\A}{\mathscr{A}}
\newcommand{\B}{\mathscr{B}} 
\newcommand{\V}{\mathscr{V}} %general labels for symm mon cat
\newcommand{\C}{\mathscr{C}} 
\newcommand{\U}{\mathscr{U}}

%%common commands (add to template)
%general
\newcommand{\FF}{\mathbb{F}} %field F
\newcommand{\NN}{\mathbb{N}} %natural number
\newcommand{\ZZ}{\mathbb{Z}} %integer
\newcommand{\QQ}{\mathbb{Q}} %rational field
\newcommand{\RR}{\mathbb{R}} %real field
\newcommand{\CC}{\mathbb{C}} %complex field
\newcommand{\KK}{\mathbb{K}} %field K
\newcommand{\GG}{\mathbb{G}} %for some functor

\newcommand{\id}{\mathrm{id}} %identity
\newcommand{\Id}{\mathrm{Id}} %big identity
\newcommand{\lcm}{\mathrm{lcm}} %least common multiple
\newcommand{\tensor}{\otimes} %tensor
\newcommand{\bigtensor}{\bigotimes} %big tensor

%algebra
\newcommand{\Gal}{\mathrm{Gal}} %galois extension
\newcommand{\Aut}{\mathrm{Aut}} %automorphism
\newcommand{\End}{\mathrm{End}} %endomorphism
\newcommand{\Hom}{\mathrm{Hom}} %set of morphism
\newcommand{\Mor}{\mathrm{Mor}} %set of morphisms (in general category, lang)
\newcommand{\Nil}{\mathrm{Nil}} %nilradical in commutatitive ring
\newcommand{\Ann}{\mathrm{Ann}} %annihilator (as ideal in commutative ring)
\newcommand{\Char}{\mathrm{char}} %characteristic
\newcommand{\coker}{\mathrm{coker}} %coker
%\let\im\undefined
\newcommand{\im}{\mathrm{im}} %image
\newcommand{\Spec}{\mathrm{Spec}} %spectrum (used in either operator theory, or spectrum as Zariski Topology)
\newcommand{\coim}{\mathrm{coim}} %coimage in homological algebra
%\newcommand{\span}{\mathrm{span}} %span in linear algebra
\newcommand{\Orb}{\mathrm{Orb}}%orbit in group theory
\newcommand{\Ob}{\mathrm{Ob}}%object in cat theory

%category theory
\newcommand{\cat}[1]{\mathsf{#1}} %general category font command
\newcommand{\Grp}{\mathsf{Grp}} %cat of group
\newcommand{\Ab}{\mathsf{Ab}} %cat of abel group
\newcommand{\Ring}{\mathsf{Ring}} %cat of ring
\newcommand{\CRing}{\mathsf{CRing}} %cat of comm ring
\newcommand{\Fld}{\mathsf{Fld}} %cat of field
\newcommand{\RMod}{\mathsf{R}\text{-}\mathsf{Mod}} %cat of R-Mod
\newcommand{\BiMod}{\mathsf{BiMod}} %cat of abel group
\newcommand{\RAlg}{\mathsf{R}\text{-}\mathsf{Alg}} %cat of R-Alg
\newcommand{\Vect}[1]{\mathsf{Vect}_\mathsf{#1}} %cat of vector spaces
\newcommand{\Top}{\mathsf{Top}} %cat of topological spaces
\newcommand{\hTop}{\mathsf{hTop}} %cat of topological spaces with homotopic classes of continuous maps
\newcommand{\Op}{\mathrm{op}}%for opposite category's notation
%\newcommand{\Sets}{\mathsf{Set}} %cat of sets
%\newcommand{\Cat}{\mathsf{Cat}} %cat of small cat
\newcommand{\Ch}[1]{\mathsf{Ch(#1)}} %cat of chain complex / cochain complex over an abelian category
\newcommand{\D}[1]{\mathsf{D(#1)}} %derived category
\newcommand{\Met}{\mathsf{Met}} %cat of metric spaces
\newcommand{\colim}{\mathrm{colim}}

%Lie Group / Lie Algebra
\newcommand{\GL}{\mathrm{GL}} %general linear
\newcommand{\SL}{\mathrm{SL}} %special linear
%\newcommand{\U}{\mathrm{U}} %unitary
\newcommand{\SO}{\mathrm{SO}} %special orthogonal
\newcommand{\SU}{\mathrm{SU}} %special unitary
\newcommand{\gl}{\mathfrak{gl}}
\let\sl\undefined
\newcommand{\sl}{\mathfrak{sl}}
\newcommand{\so}{\mathfrak{so}}
\newcommand{\su}{\mathfrak{su}}

%analysis
\newcommand{\Vol}{\mathrm{Vol}}
\newcommand{\Cont}{\C}
\newcommand{\Supp}{\mathrm{Supp}} %used for support, but can be in a lot of other use, like support in a spectrum of a commutative ring

%complex
\newcommand{\Real}{\mathrm{Re}} %real part
\newcommand{\Imag}{\mathrm{Im}} %imaginary part
%\newcommand{\Res}{\mathrm{Res}} %residue
\newcommand{\Holo}{\mathcal{O}} %represent holomorphic / analytic function
\newcommand{\Mero}{\mathcal{M}} %represent meromorphic function


%physics
\newcommand{\bR}{\mathbf{R}} %upper case position...for center of mass
\newcommand{\br}{\mathbf{r}} %position
\newcommand{\bv}{\mathbf{v}} %velocity
\newcommand{\ba}{\mathbf{a}} %acceleration
\newcommand{\bg}{\mathbf{g}} %gravity acceleration
\newcommand{\bs}{\mathbf{s}} %sometimes for extra position
\newcommand{\bF}{\mathbf{F}} %force
\newcommand{\bP}{\mathbf{P}} %momentum
\newcommand{\bp}{\mathbf{p}} %momentum operator in QM (or relativistic fourmomentum)
\newcommand{\bk}{\mathbf{k}} %wavevector
\newcommand{\bE}{\mathbf{B}} %electric field
\newcommand{\bB}{\mathbf{B}} %magnetic field
\newcommand{\bL}{\mathbf{L}} %angular momentum
\newcommand{\bS}{\mathbf{S}} %spin (QM)
\newcommand{\bJ}{\mathbf{J}} %total angular momentum
\newcommand{\bN}{\mathbf{N}} %torque
\newcommand{\bw}{\boldsymbol{\omega}} %angular velocity
\newcommand{\bOmega}{\boldsymbol{\Omega}} %big angular velocity
\newcommand{\bzero}{\mathbf{0}} %zero vector
\newcommand{\anti}{\mathrm{anti}} %antisymmetric
\newcommand{\symm}{\mathrm{symm}} %symmetric

%Probability
\newcommand{\PP}{\mathbb{P}} %probability function
\newcommand{\EE}{\mathbb{E}} %expectation value
\newcommand{\Var}{\mathrm{Var}} %use for variance
\newcommand{\Bin}{\mathrm{Bin}} %binomial distribution
\newcommand{\Ber}{\mathrm{Ber}} %bernoulli distribution
\newcommand{\Poi}{\mathrm{Poi}} %poisson distribution
\newcommand{\Unif}{\mathrm{Unif}} %uniform distribution
\newcommand{\Exp}{\mathrm{Exp}} %exponential distribution
\newcommand{\Cov}{\mathrm{Cov}} %covariance
\newcommand{\Corr}{\mathrm{Corr}} %correlation
\newcommand{\Beta}{\mathrm{Beta}} %beta distribution
\newcommand{\GammaD}{\mathrm{Gamma}} %gamma distribution
%

\bibliographystyle{plain}

%--------Meta Data: Fill in your info------
\title{Inputs for Equivariant Infinite Loop Space Machine}

\author{Zih-Yu Hsieh}

%\date{DEADLINES: Draft AUGUST 14 and Final version AUGUST 28, 2024}

\begin{document}

\begin{abstract}
Symmetric monoidal categories are central in modern algebraic topology, for instance $K$-theory relies on constructing spectra with suitable homotopy groups from spaces produced out of symmetric monoidal categories. It is known that the two inputs - pseudoalgebra over an operad $\cP$, or suitable categories described by based finite sets $\F$ - are equivalent descriptions for symmetric monoidal categories. In equivariant $K$-theory, similar machinery have been established, where the inputs for constructing $G$-spectra are described by pseudoalgebra over an operad $\cP_G$ as $G$-categories, or categories described by based finite $G$-sets $\F_G$. The equivalence of the two in the equivariant case is a natural question to ask, and associates to the question of ``what is symmetric monoidal $G$-categories''. We'll provide an exposition on the two inputs by collecting the axioms scattered throughout various pieces of the literature, while describe potential directions in equiariant equivalence.
\end{abstract}

\maketitle

\tableofcontents

\section*{Introduction}

Symmetric monoidal categories and their applications in stable homotopy theory have been developed extensively. The symmetric \emph{strict} monoidal categories (also called \emph{permutative categories}) and the $\F$-categories are central inputs for constructing spectra (where $\F$ is the standard skeleton of category of based finite sets). It is known that permutative categories are equivalent to $\cP$-algebras, where $\cP$ is the \emph{permutativity operad}. A question then arises: up to what degree are the two inputs equivalent? A more general question is answered in \cite{LMRZ26}, where the $\cP$-pseudoalgebras and \emph{strictly special $\F$-pseudoalgebras} are in fact isomorphic, showing the two inputs are directly related.

Transferring to the equivariant world, similar machines have been established: The main source of generating \emph{genuine $G$-spectra} is the $\F_G$-$G$-categories, where $\F_G$ is the category of based finite $G$-sets. Another source in the literature comes from the equivariant generalization of permutativity operad $\cP_G$, which the algebras over it is defined as \emph{permutative $G$-categories} in \cite{GM17}, while the pseudoalgebras are defined as \emph{symmetric monoidal $G$-categories} in \cite{GMMO20}. After seeing the nonequivariant equivalence, a natural followup question is whether the equivalence also holds in the equivariant case. Here we'll collect some preliminaries needed for answering the question, together with some discussions on some potential details and problems of such generalization.

\vspace{1pt}

We assume the readers have some basics about what is a category, functor, and natural transformation. In Section \ref{sec_CatPrelim} we'll discuss some tools needed for making all definitions and axioms, including 2-categories, pseudofunctors, pseudotransform, and modification between them, together with the category of small $G$-categories $\Cat_G$ we mainly work with. In Section \ref{sec_operad} we define operads over general symmetric monoidal categories, the operadic algebra, and define the nonequivariant Barratt-Eccles operad $\cP$ in $\Cat$. In Section \ref{sec_perm_G} we define the equivariant version of Barratt-Eccles operad $\cP_G$ in $\Cat_G$, together with the notion of pseudoalgebra over it. In Section \ref{sec_F_Gm} we define the category of based finite $G$-sets $\F_G$, and discuss the pseudofunctors from it to $\Cat_G$. Finally, in Section \ref{sec_results}, we discuss some potential directions (and obstacles) in generalizing the results of \cite{LMRZ26} into the equivariant case.

\vspace{1pt}

\section{Categorical Preliminaries}\label{sec_CatPrelim}
\subsection{2-Categories}\label{sub_2_cat}

\begin{defn}\label{defn_2_cat}
    A category $\C$ is a \emph{2-category} if for each object $X,Y\in\C$ (called \emph{0-cells}), the homset $\C(X,Y)$ is itself a category (called the \emph{hom-category}). Also, for all 0-cells $X,Y,Z\in\C$, the \emph{composition} $\circ:\C(Y,Z)\times\C(X,Y)\rightarrow\C(X,Z)$ is a functor.

    Every morphism $f\in \C(X,Y)$ is called a \emph{1-cell}, while given any 1-cells $f,g\in\C(X,Y)$, a morphism $\alpha:f\Rightarrow g$ in the hom-category $\C(X,Y)$ is called a \emph{2-cell}.
\end{defn}

Here, $\circ$ must be strictly unital and associative.

\begin{rem}\label{rem_bicat}
    A more general notion is that of a \emph{bicategory}, which relaxes the composition functors to be associative and unital up to coherent natural isomorphisms, and in some literature this is called 2-categories. Here we'll take the strict associativity and unitality as axioms for 2-categories. For more general notion of bicategories, cf. \cite{JY21} Section 2.1.
\end{rem}

\begin{exmp}\label{exmp_cat_of_cat}
    Let $\Cat$ denote the \emph{category of small categories}, where the 0-cells are small categories, the hom-categories $\Cat(\A,\B)$ are functors between small categories $\A$ and $\B$ (with 1-cells being functors), and the morphisms in $\Cat(\A,\B)$ are natural transformations between functors (with 2-cells being natural transformations). Thus, $\Cat$ is a 2-category.
    
    %Also, a precise description of the horizontal composition of 2-cells is as follow: Given functors $F_1,F_2\in\Cat(\A,\B)$ and $G_1,G_2\in\Cat(\B,\C)$, natural transformations $\mu:F_1\Rightarrow F_2$ in $\Cat(\A,\B)$ and $\nu:G_1\Rightarrow G_2$ in $\Cat(\B,\C)$, then the \emph{Godemen Product} $\nu*\mu := \nu \circ F\mu$ defines a horizontal composition.
\end{exmp}

\vspace{1pt}

\subsection{Pseudofunctor, Pseudotransform, and Modification}\label{sub_pseudos}

\begin{defn}[Pseudofunctor]\label{defn_pseudofunctor}
    Given two 2-categories $\A,\B$, a \emph{pseudofunctor} $F:\A\rightarrow \B$ equips with a 0-cell map, functors between hom-categories $F_{X,Y}:\A(X,Y)\rightarrow \B(F(X),F(Y))$, an invertible 2-cell $\phi_{X,Y,Z}$ between compositions
    \begin{align*}
        \xymatrix{
            \A(Y,Z)\times \A(X,Y) \ar[d]_-{\circ} \ar[rr]^-{F_{Y,Z}\times F_{X,Y}} && \B(F(Y),F(Z))\times\B(F(X),F(Y)) \ar[d]^-{\circ} \ar@{=>}[dll]_-{\phi_{X,Y,Z}}\\
            \A(X,Z) \ar[rr]_-{F_{X,Z}} && \B(F(X),F(Z))
        }
    \end{align*}
    and another invertible 2-cell $\psi_X$ between identities
    \begin{align*}
        \xymatrix{
            F(X) \ar@/^1.5pc/[rr]^{\id_{F(X)}}="1" \ar@/_1.5pc/[rr]_{F_{X,X}(\id_X)}="2"
            \ar@{=>} "1";"2"^{\psi_X}
            && F(X)
        }
    \end{align*}
    that satisfies:
    %some associativity/unitality axioms with coherence conditions. For details, cf. \cite{JY21} Section 4.1.

    %here are the results...
    \begin{itemize}
        \item[(i)] \textbf{(Associativity)} The following two pasting diagrams are equal:
        \begin{align*}
            \xymatrix{
                \A(Z,W)\times\A(Y,Z)\times\A(X,Y) \ar[d]_-{\circ\times\id} \ar[rr]^-{\prod F_{\_,\_}} && \B(F(Z),F(W))\times \B(F(Y),F(Z)) \times \B(F(X),F(Y)) \ar[d]^-{\circ\times \id} \ar@{=>}[dll]_-{\phi_{Y,Z,W}\times\id}\\
                \A(Y,W)\times\A(X,Y) \ar[d]_-{\circ} \ar[rr]^-{F_{Y,W}\times F_{X,Y}} && \B(F(Y),F(W))\times \B(F(X),F(Y)) \ar[d]^-{\circ} \ar@{=>}[dll]_-{\phi_{X,Y,W}}\\
                \A(X,W)\ar[rr]_-{F_{X,W}} && \B(F(X),F(W))
            }
        \end{align*}
        equals
        \begin{align*}
            \xymatrix{
                \A(Z,W)\times\A(Y,Z)\times\A(X,Y) \ar[d]_-{\id\times\circ} \ar[rr]^-{\prod F_{\_,\_}} && \B(F(Z),F(W))\times \B(F(Y),F(Z)) \times \B(F(X),F(Y)) \ar[d]^-{\id\times\circ} \ar@{=>}[dll]_-{\id\times\phi_{X,Y,Z}}\\
                \A(Z,W)\times\A(X,Z) \ar[d]_-{\circ} \ar[rr]^-{F_{Z,W}\times F_{X,Z}} && \B(F(Z),F(W))\times \B(F(X),F(Z)) \ar[d]^-{\circ} \ar@{=>}[dll]_-{\phi_{X,Z,W}}\\
                \A(X,W)\ar[rr]_-{F_{X,W}} && \B(F(X),F(W))
            }
        \end{align*}
        \item[(ii)] \textbf{(Unitality)} The following two unit diagrams commute:
        \begin{align*}
            \xymatrix{
                \A(X,Y)\times\A(X,X)\ar[rr]^-{F_{X,Y}\times F_{X,X}} && \B(F(X),F(Y))\times \B(F(X),F(X)) \ar[dd]^-{\circ} \ar@{=>}[ddll]_-{\id \mathrm{*}\psi_X}\\
                \A(X,Y)\times \mathrm{*} \ar[u]^-{\id\times \id_X} \ar[d]_-{F_{X,Y}\times \id_{F(X)}}\\
                \B(F(X),F(Y))\times \B(F(X),F(X)) \ar[rr]_-{\circ} && \B(F(X),F(Y))
            }
        \end{align*}
        \begin{align*}
            \xymatrix{
                \A(Y,Y)\times\A(X,Y)\ar[rr]^-{F_{Y,Y}\times F_{X,Y}} && \B(F(Y),F(Y))\times \B(F(X),F(Y)) \ar[dd]^-{\circ} \ar@{=>}[ddll]_-{\psi_Y\mathrm{*}\id}\\
                \mathrm{*}\times\A(X,Y) \ar[u]^-{\id_Y\times \id} \ar[d]_-{\id_{F(Y)}\times F_{X,Y}}\\
                \B(F(Y),F(Y))\times \B(F(X),F(Y)) \ar[rr]_-{\circ} && \B(F(X),F(Y))
            }
        \end{align*}
        where $\id_X:*\rightarrow\A(X,X)$ is the functor indicating identity of $X$.
    \end{itemize}

    If $\psi_X$ is the identity 2-cell for all 0-cell $X$, then the pseudofunctor $F$ is called \emph{normal}.
\end{defn}

\vspace{1pt}

\begin{defn}[Pseudotransform]\label{defn_pseudotransform}
    Fix 2-categories $\A,\B$, let $(F,\phi,\psi), (G,\varphi,\vartheta):\A\rightarrow\B$ be two pseudofunctors. A \emph{pseudotransform} $\mu:(F,\phi,\psi)\rightarrow (G,\varphi,\vartheta)$ collects 1-cells $\mu_X:F(X)\rightarrow G(X)$ for each 0-cell $X\in\A$, and for each 1-cell $f:X\rightarrow Y$ in $\A$, associates a 2-cell $\mu f$ such that 
    \begin{align*}
        \xymatrix{
            F(X) \ar[r]^-{\mu_X} \ar[d]_-{Ff} & G(X) \ar[d]^-{Gf} \ar@{=>}[dl]_-{\mu f}\\
            F(Y) \ar[r]_-{\mu_Y} & G(Y)
        }
    \end{align*}
    commutes. They also satisfy:
    %coherent associativity/unitality axioms, with details in \cite{JY21} Section 4.2.
    %details here
    \begin{itemize}
        \item[(i)] (\textbf{Associativity}) Let $f:X\rightarrow Y$, $g:Y\rightarrow Z$ be 1-cells in $\A$, the following two pasting diagrams are equal:
        \begin{align*}
            \xymatrix{
                F(X) \ar[dd]_-{F(g\circ f)} \ar[dr]^-{Ff} \ar[rr]^-{\mu_X} && G(X) \ar@{=>}[dl]_-{\mu f} \ar[dr]^-{Gf}\\
                & F(Y) \ar[rr]^-{\mu_Y} \ar@{=>}[l]_-{\phi_{g,f}} \ar[dl]^-{Fg} && G(Y) \ar[dl]^-{Gg} \ar@{=>}[dlll]_-{\mu g}\\
                F(Z) \ar[rr]_-{\mu_Z} && G(Z)
            } \quad = \quad \xymatrix{
                F(X) \ar[dd]_-{F(g\circ f)} \ar[rr]^-{\mu_X} && G(X) \ar[dr]^-{Gf} \ar[dd]_-{G(g\circ f)} \ar@{=>}[ddll]_-{\mu(g\circ f)}\\
                &&& G(Y) \ar[dl]^-{Gg} \ar@{=>}[l]_-{\varphi_{g,f}}\\
                F(Z) \ar[rr]_-{\mu_Z} && G(Z)
            }
        \end{align*}
        \item[(ii)] (\textbf{Unitality}) The following pasting diagrams are equal:
        \begin{align*}
            \xymatrix{
                F(X) \ar@(dl,ul)[dd]_-{F(\id_X)} \ar@(dr,ur)[dd]^-{\id_{F(X)}} \ar@{=>}[dd]_-{\psi_X} \ar[rr]^-{\mu_X} && G(X) \ar[dd]^-{\id_{G(X)}}\\
                \\
                F(X) \ar[rr]_-{\mu_X} && G(X)
            }\quad = \quad \xymatrix{
                F(X) \ar[dd]_-{F(\id_X)} \ar[rr]^-{\mu_X} && G(X) \ar@{=>}[ddll]_-{\mu(\id_X)} \ar@(dl,ul)[dd]_-{G(\id_X)} \ar@(dr,ur)[dd]^-{\id_{G(X)}} \ar@{=>}[dd]^-{\vartheta_X}\\
                \\
                F(X) \ar[rr]_-{\mu_X} && G(X)
            }
        \end{align*}
    \end{itemize}
\end{defn}

\vspace{1pt}

\begin{defn}[Modification]\label{defn_modification}
    Continue with the same notation as above. Let $\mu,\nu:(F,\phi,\psi)\rightarrow (G,\varphi,\vartheta)$ be two pseudotransforms, a \emph{modification} $\lambda:\mu\Rightarrow\nu$ is a collection of 2-cells $\lambda_X:\mu_X\Rightarrow\nu_X$ for each 0-cell $X\in\A$, such that the following two pasting diagrams are equal, for each 1-cell $f:X\rightarrow Y$ in $\A$:
    \begin{align*}
        \xymatrix{
            F(X) \ar[dd]_-{Ff} \ar@(ur,ul)[rr]^-{\mu_X} \ar@(dr,dl)[rr]_-{\nu_X} \ar@{=>}[rr]^-{\lambda_X} && G(X) \ar@{=>}[ddll]^-{\nu f} \ar[dd]^-{Gf}\\
            \\
            F(Y) \ar[rr]_-{\nu_Y} && G(Y)
        }\quad = \quad \xymatrix{
            F(X) \ar[dd]_-{Ff} \ar[rr]^-{\mu_X} && G(X) \ar[dd]^-{Gf} \ar@{=>}[ddll]_-{\mu f}\\
            \\
            F(Y) \ar@(ur,ul)[rr]^-{\mu_Y} \ar@{=>}[rr]^-{\lambda_Y} \ar@(dr,dl)[rr]_-{\nu_Y} && G(Y)
        }
    \end{align*}

\end{defn}

\vspace{1pt}

%There are obvious composition laws between pseudofunctors, the pseudotransforms, and also the modifications. (\textbf{Depending on situation, fill in details, or delete this line})

%\vspace{1pt}

%If enlarging the notion of functors and natural transformations to pseudofunctors and pseudotransforms, $\Cat$ will be enlarged into another 2-categories. However, the detail of composition and its unitality/associativity is omitted here for space limit.

\subsection{$G$-Categories}\label{sub_G_cat}

\begin{defn}\label{defn_G_cat}
    A small category $\A$ is a \emph{$G$-category} if the object and morphism sets $\Ob(\A), \Mor(\A)$ are both (left) $G$-sets, such that the following four structure maps $(S,T,I,C)$ (stands for \emph{source, target, identity, and composition}) are all $G$-equivariant:
    \begin{itemize}
        \item $S,T:\Mor(\A)\rightarrow\Ob(\A)$ by $S(X\xrightarrow{f}Y):= X$, and $T(X\xrightarrow{f}Y):= Y$.
        \item $I:\Ob(\A)\rightarrow\Mor(\A)$ by $I(X):= \id_X$.
        \item $C:\Mor(\A)\times_{S,T}\Mor(\A)\rightarrow\Mor(\A)$ by $C(g,f):= g\circ f$.
    \end{itemize}
    Here, $\Mor(\A)\times_{S,T}\Mor(\A):= \{(g,f)\in \Mor(\A)^2\ |\ S(g)=T(f)\}$, precisely the pair of morphisms that're composable, and it endows $G$-set structure via restricting the diagonal $G$-action on $\Mor(\A)^2$. This is well-defined due to the assumption $S,T$ are $G$-maps.
\end{defn}

Another equivalent characterization can be viewed as a group homomorphism $G\rightarrow\Aut_{\Cat}(\A)$, which assigns each $g\in G$ an automorphism functor on $\A$ such that composition respects $G$'s multiplication structure, and identity gets sent to the identity functor. So, from now on we view each element $g\in G$ as a functor, whenever a category is specified as a $G$-category. 

\vspace{1pt}

\begin{defn}\label{defn_cat_of_G_cat}
Let $\Cat_G$ denotes the 2-category of small $G$-categories, with (nonequivariant) functors as $1$-cells, natural transformations as $2$-cells. The hom-category between any two $G$-categories $\Cat_G(\A,\B)$ also has a natural $G$-action via conjugation.  Which, the $G$-fixed subcategory of the hom-category $[\Cat_G(\A,\B)]^G$ precisely describes all $G$-functors, or the functors $F:\A\rightarrow\B$ such that $g\circ F=F\circ g$ for all $g\in G$.
\end{defn}

\vspace{1pt}

Another important notion is the twisted $G$-action on products when given a group homomorphism $\alpha:G\rightarrow\Sigma_n$. Given a category $\A$, the $n$-fold product $\A^\mathbf{n}$ has a natural left $\Sigma_n$-action via permuting the coordinates; if $\A$ is a $G$-category, $\A^\mathbf{n}$ also endows a natural diagonal $G$-action. These give a natural left $(G\times \Sigma_n)$-action on $\A^\mathbf{n}$.

\begin{defn}
    Let $\A^{\mathbf{n}^\alpha}$ be a $G$-category with the underlying category being $\A^\mathbf{n}$, with the $G$-action defined as follows:
    \begin{align*}
        g\cdot_\alpha (X_1,...,X_n):= (g\cdot X_{\alpha(g)^{-1}(1)},...,g\cdot X_{\alpha(g)^{-1}(n)})
    \end{align*}
\end{defn}
Another characterization can be given as follow: Given $\alpha$'s graph subgroup $\Gamma_\alpha:= \{(g,\alpha(g))\in G\times\Sigma_n\}_{g\in G}$, the projection $\pi_G:G\times\Sigma_n\rightarrow G$ restricts to an isomorphism $\Gamma_\alpha\xrightarrow{\sim}G$. The $\alpha$-twisted action is induced by the pullback along this isomorphism.

\vspace{1pt}

\subsection{Chaotic-Forgetful Adjunction}\label{sub_chaot_adj}

Given any small category $\A\in\Cat$, its object $\Ob(\A)$ forms a set; moreover, any functor $F:\A\rightarrow\B$ also has an underlying map of objects $F:\Ob(\A)\rightarrow\Ob(\B)$. Conversely, given a set, it has certain category structure:

\begin{defn}\label{defn_chaot}
    Given a set $X$, define $\mathscr{E}X$ to be the \emph{chaotic category} with underlying object set $\Ob(\mathscr{E}X):=X$, by encoding a unique isomorphism $x\xrightarrow{\sim} y$ between any element $x,y\in X$. The composition 
    \begin{align}
        (x\xrightarrow{\sim}y \xrightarrow{\sim}z) = (x\xrightarrow{\sim}z)
    \end{align}
    using the uniqueness.

    \vspace{1pt}

    Given a set map $f:X\rightarrow Y$, define the functor $\mathscr{E}f:\mathscr{E}X\rightarrow\mathscr{E}Y$ so that the underlying object map is $f$, and for each morphism $x\xrightarrow{\sim}y$ in $\mathscr{E}X$, $\mathscr{E}f(x\rightarrow y):= f(x)\rightarrow f(y)$, the unique morphism between $f(x),f(y)$ in $\mathscr{E}Y$.
\end{defn}

Inspecting the definition, we see $\mathscr{E}:\Sets\rightarrow\Cat$ is a well-defined functor, called the \emph{chaotic functor}.

Notice that if a functor $F:\A\rightarrow\B$ has a chaotic target $\B$, then all the morphisms $f:X\rightarrow Y$ in $\A$ must be sent to the unique morphism $F(X)\rightarrow F(Y)$ in $\B$, hence $F$ is dictated by the object map $F:\Ob(\A)\rightarrow\Ob(\B)$. Such uniqueness specifies an adjunction $(\Ob\dashv \mathscr{E})$.

The consequence is the following natural isomorphism of homset for any two sets $X,Y$:
\begin{align}
    \Sets(X,Y)=\Sets(\Ob(\mathscr{E}X),Y)\cong \Cat(\mathscr{E}X,\mathscr{E}Y)
\end{align}

\begin{rem}\label{rem_equiv_chaot}
    Let $\Sets_G$ denotes the category of (left) $G$-sets, with all maps between $G$-sets (not necessarily equivariant) being the morphisms, the above adjunction can be extended to $\Ob:\Cat_G\rightarrow \Sets_G$, and $\mathscr{E}:\Sets_G\rightarrow\Cat_G$. Similar notions hold when transfer from left to right $G$-sets/categories.
\end{rem}

In this exposition we only need $\mathscr{E}$ to be a right adjoint, there are richer details demonstrated in \cite{GM17} Section 3.

\section{Operads, Operadic Algebra, and Operad $\cP$}\label{sec_operad}
\subsection{Symmetric Monoidal Categories}\label{sub_symmon}

\begin{defn}\label{defn_symmon}
    A category $\mathcal{V}$ is a \emph{symmetric monoidal category}, if there is a bifunctor $\tensor:\mathcal{V}\times\mathcal{V}\rightarrow\mathcal{V}$ (called \emph{monoidal product}), and an object $\kappa\in\mathcal{V}$ (called \emph{unit/monoidal unit}), together with the following natural isomorphisms:
    \begin{itemize}
        \item The \emph{associator} $a:(\_ _1\tensor \_ _2)\tensor \_ _3 \xrightarrow{\sim} \_ _1\tensor(\_ _2\tensor \_ _3)$.
        \item The \emph{left unit} $\lambda:\kappa\tensor\_\xrightarrow{\sim} \Id_{\mathcal{V}}$ and the \emph{right unit} $\rho:\_\tensor \kappa\xrightarrow{\sim}\Id_\mathcal{V}$.
        \item The \emph{symmetric braiding} $B:\_ _1\tensor\_ _2\xrightarrow{\sim}\_ _2\tensor\_ _1$, with extra property $B_{Y,X}:Y\tensor X\rightarrow X\tensor Y$ equals to $B_{X,Y}^{-1}:X\tensor Y\rightarrow Y\tensor X$.
    \end{itemize} 

    They also satisfy the following axioms:
    \begin{itemize}
        \item \textbf{(Triangle)}
        \[\xymatrix{
            (X\tensor \kappa)\tensor Y \ar[dr]_-{\rho_X \tensor \id_Y} \ar[rr]^-{a_{X,\kappa,Y}} && X\tensor (\kappa\tensor Y) \ar[dl]^-{\id_X\tensor \lambda_Y}\\
            & X\tensor Y
        }\]
        \item \textbf{(Pentagon)}
        \[\xymatrix{
            ((X\tensor Y)\tensor Z)\tensor W \ar[rr]^-{a_{X\tensor Y,Z,W}} \ar[d]_-{a_{X,Y,Z}\tensor \id_W} && (X\tensor Y)\tensor(Z\tensor W) \ar[d]^-{a_{X,Y,Z\tensor W}}\\
            (X\tensor (Y\tensor Z))\tensor W \ar[dr]_-{a_{X,Y\tensor Z,W}} && X\tensor (Y\tensor (Z\tensor W)) \\
            & X\tensor ((Y\tensor Z)\tensor W) \ar[ur]_-{\id_X\tensor a_{Y,Z,W}}
        }\]
        \item \textbf{(Hexagon)}
        \[\xymatrix{
            (X\tensor Y)\tensor Z \ar[r]^-{a_{X,Y,Z}} \ar[d]_-{B_{X,Y}\tensor \Id_Z} & X\tensor (Y\tensor Z)\ar[r]^-{B_{X,Y\tensor Z}} & (Y\tensor Z)\tensor X \ar[d]^-{a_{Y,Z,X}}\\
            (Y\tensor X)\tensor Z \ar[r]_-{a_{Y,X,Z}} & Y\tensor (X\tensor Z) \ar[r]_-{\id_Y\tensor B_{X,Z}} & Y\tensor (Z\tensor X)
        }\]
        %\[\xymatrix{
        %    X\tensor (Y\tensor Z) \ar[d]_-{\id_X\tensor B_{Y,Z}} & (X\tensor Y)\tensor Z \ar%[l]_-{a_{X,Y,Z}} \ar[r]^-{B_{X\tensor Y,Z}} & Z\tensor (X\tensor Y)\\
        %    X\tensor (Z\tensor Y) & (X\tensor Z)\tensor Y \ar[l]^-{a_{X,Z,Y}} \ar[r]_-{B_{X,Z}\tensor \id_Y} & (Z\tensor X)\tensor Y \ar[u]_-{a_{Z,X,Y}}
        %}\]
    \end{itemize}
\end{defn}
Based on \emph{Maclane's Coherence Theorem}, the Triangle and Pentagon axioms guarantees all finite sequences of $a,\lambda,\rho$ and their inverses agree, namely there is a unique isomorphism $P_1\rightarrow P_2$ consists of compositions of $a,\lambda,\rho$, where $P_i$ are (possibly) different orders of performing product on $X_1,...,X_n\in\mathcal{V}$. For details, cf. \cite{M98} Chapter XI.

The category of sets $\Sets$, small categories $\Cat$ in Example \ref{exmp_cat_of_cat}, and small $G$-categories $\Cat_G$ in Definition \ref{defn_cat_of_G_cat} are central examples for this, where the cartesian product $\A\times\B$ serves as monoidal product, and a terminal set/category $*$ serves as units for both.

\vspace{1pt}

\subsection{Operads, Operadic Algebras}\label{sub_operad}

\begin{defn}\label{defn_operad}
Given symmetric monoidal category $\mathcal{V}$, a \emph{$\Sigma$-operad} $\C$ in $\mathcal{V}$ is a collection of objects $\{\C(j)\in\mathcal{V}\}_{j\in\NN}$, each $\C(j)$ equipsx with a right $\Sigma_j$-action (a group homomorphism $\Sigma_j^\Op\rightarrow \Aut_\mathcal{V}(\C(j))$), a \emph{unit map} $\eta:1\rightarrow \C(1)$, together with composition maps
\[\gamma:\C(k)\tensor \C(j_1)\tensor...\tensor \C(j_k)\rightarrow \C(j_1+...+j_k)\]
that satisfies the following axioms:
\begin{itemize}
    \item \textbf{(Unit)}
    \[\xymatrix{
        \C(j)\tensor \kappa^{\tensor j} \ar[dr]_-{\id_{\C(j)}\tensor \eta^{\tensor j}} \ar[rr]^-{\sim} && \C(j)\\
        & \C(j)\tensor \C(1)^{\tensor j} \ar[ur]_-{\gamma}
    },\quad\quad \xymatrix{
        \kappa\tensor \C(j) \ar[dr]_-{\eta\tensor \id_{\C(j)}} \ar[rr]^-{\sim} && \C(j)\\
        & \C(1)\tensor \C(j) \ar[ur]_-{\gamma}
    }\]
    \item \textbf{(Associativity)} Fix natural numbers $k$, $j = \sum_{s=1}^k j_s$, define index $g_s:= j_1+...+j_s$ for $1\leq s\leq k$. For each $i_1,...,i_{g_k}=i_j$, with sum $i=\sum_{r=1}^j i_r$, and $h_s:= i_{g_{(s-1)}+1}+...+i_{g_s}$ (so $i=\sum_{s=1}^k h_s$ also), the following diagram commutes:
    \[\xymatrix{
        \left[\C(k)\tensor \left(\bigtensor_{s=1}^k \C(j_s)\right)\right]\tensor \left(\bigtensor_{r=1}^j \C(i_r)\right) \ar[dd]_-{\text{shuffle}} \ar[rr]^-{\gamma\tensor \id} && \C(j) \tensor \left(\bigtensor_{r=1}^j \C(i_r)\right) \ar[d]^-{\gamma}\\
        && \C(i)\\
        \C(k) \tensor \left[\bigtensor_{s=1}^k \left(\C(j_s)\tensor \left(\bigtensor_{q=1}^{j_s} \C(i_{g_{(s-1)}+q})\right)\right)\right] \ar[rr]_-{\id_{\C(k)} \tensor (\bigtensor_s \gamma)} && \C(k)\tensor \left(\bigtensor_{s=1}^k \C(h_s)\right) \ar[u]_-{\gamma}
    }\]
    \item \textbf{(Equivariance)}
    \begin{itemize}
        \item Given $\sigma\in\Sigma_k$, $j_1+...+j_k = \Sigma_j$, let $\sigma(j_1,...,j_k)\in\Sigma_j$ denotes ``permutation of $k$-blocks of letters'', by $(j_1,...,j_k)\mapsto (j_{\sigma^{-1}(1)},...,j_{\sigma^{-1}(k)})$.
        Then, the following diagram commutes:
        \[\xymatrix{
            \C(k)\tensor \left(\bigtensor_{r=1}^k \C(j_r)\right) \ar[d]_-{\gamma} \ar[rr]^-{\sigma\tensor \sigma^{-1}} && \C(k)\tensor \left(\bigtensor_{r=1}^k \C(j_{\sigma(r)})\right) \ar[d]^-{\gamma}\\
            \C(j) \ar[rr]_-{\sigma(j_{\sigma(1)},...,j_{\sigma(k)})} && \C(j)
        }\]
        \item Given $\tau_r\in \Sigma_{j_r}$, denote $\tau_1\oplus ... \oplus \tau_k\in \Sigma_j$ by letting each $\tau_r$ permute the $j_r$th block in $j$. 
        %(for instance, let $\tau_1$ acts on $1,...,j_1$; let $\tau_2$ acts on $j_1+1,...,j_1+j_2$, etc.). 
        Then, the following diagram commutes:
        \[\xymatrix{
            \C(k)\tensor \left(\bigtensor_{r=1}^k \C(j_r)\right)\ar[d]_-{\gamma} \ar[rr]^-{\id_{\C(k)}\tensor (\bigtensor_r \tau_r)} && \C(k)\tensor \left(\bigtensor_{r=1}^k \C(j_r)\right) \ar[d]^-{\gamma}\\
            \C(j) \ar[rr]_{\tau_1\oplus...\oplus\tau_k} && \C(j)
        }\]
    \end{itemize}
\end{itemize}

If $\C(0)=\kappa$ (with $\gamma:\C(0)\rightarrow \C$ being $\id_{\C(0)}$), the operad is \emph{reduced}.
\end{defn}

\begin{defn}\label{defn_operad_alg}
    Given an operad $\C$ over a symmetric monoidal category $\mathcal{V}$, a \emph{$\C$-algebra} is an object $A\in\mathcal{V}$, with collections of morphisms $\theta:\C(j)\tensor A^j\rightarrow A$, that satisfies the following axioms:
    \begin{itemize}
        \item \textbf{(Associativity)} Given $j = \sum_{i=1}^k j_i$, the following diagram commutes:
        \[\xymatrix{
            \left[\C(k)\tensor \left(\bigtensor_{i=1}^k \C(j_i)\right)\right] \tensor A^{\tensor j} \ar[dd]_-{\text{shuffle}} \ar[rr]^-{\gamma\tensor \id_{A^{\tensor j}}} && \C(j)\tensor A^{\tensor j} \ar[d]^-{\theta}\\
            & & A\\
            \C(k)\tensor \left[\bigtensor_{i=1}^k \left(\C(j_i)\tensor A^{\tensor j_i}\right)\right] \ar[rr]_-{\id_{\C(k)}\tensor \theta^k} && \C(k)\tensor A^{\tensor k} \ar[u]_-{\theta}
        }\]
        (Note: This part uses the isomorphism $A^{\tensor j}\cong\bigtensor_{i=1}^k A^{\tensor j_i})$.
        \item \textbf{(Unit)} The following diagram commutes:
        \[\xymatrix{
            \kappa \tensor A \ar[dr]_-{\eta\tensor \id_A} \ar[rr]^-{\sim} && A\\
            & \C(1)\tensor A \ar[ur]_-{\theta}
        }\]
        \item \textbf{(Equivariance)} Given any permutation $\sigma\in\Sigma_j$, the following diagram commutes:
        \[\xymatrix{
            \C(j)\tensor A^{\tensor j} \ar[dr]_-{\theta} \ar[rr]^-{\sigma\tensor \sigma^{-1}} && \C(j)\tensor A^{\tensor j} \ar[dl]^-{\theta}\\
            & A
        }\]
    \end{itemize}
\end{defn}

\begin{exmp}
    The \emph{associativity operad} $\M$ in $\Sets$ is defined as $\M(j):=\Sigma_j$ (the symmetric group for cardinality $j$), which endows a natural right $\Sigma_j$-action via right multiplication. The operadic structure map is given by 
    \begin{align}
        \gamma:\Sigma_k\times\prod_{r=1}^k\Sigma_{j_r}\rightarrow\Sigma_j,\quad \gamma(\sigma;\tau_1,...,\tau_k):= \tau_{\sigma^{-1}(1)}\oplus...\oplus\tau_{\sigma^{-1}(k)}\circ \sigma(j_{\sigma(1)},...,j_{\sigma(k)})
    \end{align}
    where $j=\sum_r j_r$, $\sigma(j_1,...,j_k)\in \Sigma_j$ denotes the block permutation of $j_1,...,j_k$ within $j$ (in order), and $\tau_1\oplus ... \tau_{k}$ is the image of the embedding $\oplus:\prod_{r=1}^k\Sigma_{j_r}\rightarrow\Sigma_j$, by viewing the $r$th block of $\mathbf{j}$ (with $j_r$-elements) as finite set $\mathbf{j}_r$. Inspecting the definitions, the algebra over $\M$ precisely describes the set-theoretic monoids. For more details, cf.  \cite{GM17} Section 4.1.
\end{exmp}

\subsection{Permutativity Operad $\cP$} \label{sub_perm_operad}
Now, recall that the chaotic functor $\mathscr{E}:\Sets\rightarrow\Cat$ is a right adjoint of $\Ob:\Cat\rightarrow\Sets$, in particular it preserves products. Hence, the chaotic categorification $\cP(j):= \mathscr{E}\M(j)$ also has a natural operadic structure, via the identification
\begin{align}
    \mathscr{E}\M(k)\times \prod_{r=1}^k \mathscr{E}\M(j_r)\cong \mathscr{E}\left(\M(k)\times\prod_{r=1}^k\M(j_r)\right) \xrightarrow{\mathscr{E}\gamma}\mathscr{E}\M(j).
\end{align}
This is called the \emph{permutativity operad in $\Cat$}, or the \emph{Barratt-Eccles operad}. As an abuse of notation, the operadic structure map of $\cP$ is also denoted $\gamma$.

\section{$G$-Operad $\cP_G$ and its Pseudoalgebra}\label{sec_perm_G}
\subsection{The $G$-Operad $\cP_G$}\label{sub_perm_G}
First we consider the permutativity operad $\cP$ in $\Cat$. If we view each $\cP(j)$ as a $G$-trivial $G$-category, it is also an operad in $\Cat_G$. Its underlying algebra is called the \emph{naive permutative $G$-category}, as specified in \cite{GM17} Section 4.1, it produces a naive $G$-spectrum under equivariant infinite loop space theory. In \cite{GM17} Section 4.2, a richer equivariant structure is introduced which produces the \emph{equivariant Barratt Eccles $G$-operad} $\cP_G$ as follows:
\begin{defn}\label{defn_perm_G}
    The $j$th category $\cP_G(j) := \Cat(\mathscr{E}G,\cP(j))$, with the operadic structure map determined by the fact that $\Cat(\mathscr{E}G,\_)$ preserves product:
    \begin{align}
        \Cat(\mathscr{E}G,\cP(k))\times \prod_{r=1}^k\Cat(\mathscr{E}G,\cP(j_r))\cong \Cat\left(\mathscr{E}G,\cP(k)\times \prod_{r=1}^k\cP(j_r)\right)\xrightarrow{\gamma\circ\_}\Cat(\mathscr{E}G,\cP(j))
    \end{align}
    By abuse of notation, we'll denote the operadic structure map as $\gamma$ again.
\end{defn}

Another formulation can be done by defining $P_G(j):= \Sets(G,\Sigma_j)$ in $\Sets$, then utilize chaotic functor promoting this to a chaotic operad $\mathscr{E}\Sets(G,\Sigma_j)\cong \cP_G(j)$. Either way there is a right $G$-action on each $\cP_G(j)$, via right multiplying the inverse before plugging in the function:
\begin{align}
    (f:G\rightarrow\Sigma_j)\mapsto (f(\_\cdot g):G\rightarrow\Sigma_j)
\end{align}

\subsection{The 2-Category $\cP_G$-$\PsAlg$}\label{sub_P_G_psalg}

The general notion of operadic pseudoalgebra is defined in \cite{CG14}, \cite{GMMO20}, \cite{LMRZ26}, and \cite{Y24 }. Here we only take the case when the operad is $\cP_G$, together with some description of $G$-equivariance.

\begin{defn}[$\cP_G$-pseudoalgebra]\label{defn_P_G_psalg}
    A $\cP_G$-pseudoalgebra is a $G$-category $\A$ and operad action functors
    \begin{align*}
        \theta_j:\cP_G(j)\times \A^j\rightarrow\A
    \end{align*}
    for each $n\in\NN$, together with an invertible natural transformation $\phi_1$ in the diagram
    \begin{align}
        \xymatrix{ 
        \A\ar[r]^{\iota}\ar[dr]_{\Id_\A} & \cP_G(1)\times\A\ar[d]^{\theta_1}\ar@{=>}[dl]+<26pt>^>>>>{\phi_1}\\ 
        & C
        }
    \end{align}
    An invertible natural transformation $\phi=\phi(k;j_1,...,j_k)$ for each diagram 
    \begin{align}
        \xymatrix{
            \cP_G(k) \times \left(\prod_{r=1}^k \cP_G(j_r)\times \A^{j_r}\right) \ar[dd]_-{\text{shuffle}} \ar[rr]^-{\id\times\prod_r\theta_{j_r}} && \cP_G(k)\times\A^k \ar[d]^-{\theta_k} \ar@{=>}[ddll]^-{\phi}\\
            && \A\\
            \cP_G(k)\times \prod_{r=1}^k\cP_G(j_r)\times \A^{j_+} \ar[rr]_-{\gamma\times\id} && \cP_G(j_+)\times\A^{j_+} \ar[u]_-{\theta_{j_+}}
        }
    \end{align}
    such that the following strict equivariance conditions hold for $\theta_j$ and $\phi$:
    \begin{itemize}
        \item[(i)]\textbf{(Equivariance of $\theta_j$)} The following two diagrams commutes for $\sigma\in\Sigma_j$ and $g\in G$:
        \begin{align}
            \xymatrix{
                \cP_G(j)\times \A^j \ar[dr]_-{\theta_j}  \ar[rr]^-{\sigma\times\sigma^{-1}} && \cP_G(j)\times \A^j \ar[dl]^-{\theta_j}\\
                & \A
            },\quad\quad \xymatrix{
                \cP_G(j)\times \A^j \ar[d]_-{\theta_j}  \ar[rr]^-{g\times g^{-1}} && \cP_G(j)\times \A^j \ar[d]^-{\theta_j}\\
                \A \ar[rr]_-{g^{-1}} && \A
            }
        \end{align}
        The left is the operadic equivariance, the right is $G$-equivariance by equipping $\cP(j)\times\A^j$ with the product $G$-action.
        %\textbf{I think if assuming the commutation for $G$ also, the case $j=1$ implies $\theta_j$ collapses all $G$-action on $\A$ to be trivial}
        \item[(ii)]\textbf{(Equivariance of $\phi$)} The following equalities hold when given $\sigma\in\Sigma_k$, and $\tau_r\in\Sigma_{j_r}$: 
        \begin{align*}
            \phi(k;j_1,...,j_k) = \phi(k;j_{\sigma(1)},...,j_{\sigma_k}) \circ (\sigma\times\sigma^{-1})
        \end{align*}
        \begin{align*}
            \phi(k;j_1,...,j_k)=\phi(k;j_1,...,j_k)\circ (\id\times \prod_r(\tau_r\times\tau_r^{-1}))
        \end{align*}
        %\textbf{(There are subtleties in $G$-equivariance)}
    \end{itemize}
    The transformations also need to satisfy the following coherence properties:
    \begin{itemize}
        \item[(iii)] \textbf{(Operadic Composition)} The following pasting diagrams are equal, given $j_+=\sum j_r$, $i_r = \sum_si_{r,s}$, and $i_+ = \sum_{r,s}i_{r,s}$ with $1\leq r\leq k$ and $1\leq s\leq j_r$:
        \item[] \begin{align*}
            \xymatrix{
                \cP_G(k)\times \prod_r\left(\cP_G(j_r)\times \prod_s\left(\cP_G(i_{r,s})\times \A^{i_{r,s}}\right)\right) \ar[r]^-{\id\times \prod(\id\times \theta_{i_{r,s}})} \ar[dd]_{\gamma\times\id} & Q_G(k)\times \prod_r\left(\cP_G(j_r)\times \A^{j_r}\right) \ar[dd]_-{\gamma} \ar[dr]^-{\id\times \prod_{\theta_{j_r}}}\\
            && \cP_G(k)\times \A^k \ar[dd]^-{\theta_k} \ar@{=>}[dl]^-{\phi}\\
            \cP_G(j_+)\times \prod_{r,s}\left(\cP_G(i_{r,s})\times \A^{i_{r,s}}\right) \ar[r]^-{\id\times \prod \theta_{i_{r,s}}} \ar[dr]_-{\gamma\times \id} & \cP_G(j_+)\times \A^{j_+} \ar[dr]^-{\theta_{j_+}} \ar@{=>}[d]^-{\phi}\\
            & \cP_G(i_+)\times \A^{i_+} \ar[r]_-{\theta_{j_+}} & \A
            }
        \end{align*}
        equals
        \begin{align*}
            \xymatrix{
                \cP_G(k)\times \prod_r\left(\cP_G(j_r)\times \prod_s\left(\cP_G(i_{r,s})\times\A^{i_{r,s}}\right)\right) \ar[r]^-{\id\times \prod(\id\times \prod \theta_{i_{r,s}})} \ar[dd]_-{\gamma\times\id} \ar[dr]_-{\id\times \prod(\gamma\times\id)} & \cP_G(k)\times \prod_r\left(\cP_G(j_r)\times \A^{j_r}\right) \ar[dr]^-{\id\times \prod\theta_{j_r}} \ar@{=>}[d]^-{\id\times \prod\phi}\\
                & \cP_G(k)\times \prod_r\left(\cP_G(i_r)\times \A^{i_r}\right) \ar[r]_-{\id\times\prod\gamma} \ar[dd]_-{\gamma\times\id} & \cP_G(k)\times\A^k \ar[dd]^-{\theta_k} \ar@{=>}[ddl]^-{\phi}\\
                \cP_G(j_+)\times\prod_{r,s}\left(\cP_G(i_{r,s})\times \A^{i_{r,s}}\right) \ar[dr]_-{\gamma\times\id}\\
                & \cP_G(i_+)\times\A^{i_+} \ar[r]_-{\theta_{i_+}} & \A
            }
        \end{align*}
        The commuting square above is due to the associativity of operad.
        \item[(iv)]\textbf{(Operadic Identity)} The following two pairs of pasting diagrams are equal:
        \begin{align*}
            \xymatrix{
                \cP_G(k)\times \A^k\ar[r]^-{\id\times\iota^k} & \cP_G(k)\times \left(\cP_G(1)\times\A\right)^k \ar[r]^-{\id\times \theta_1^k} \ar[d]_-{\text{shuffle}} & \cP_G(k)\times \A^k \ar[dd]^-{\theta_k} \ar@{=>}[ddl]^-{\phi}\\
                & \cP_G(k)\times \cP_G(1)^k\times \A^k \ar[d]_-{\gamma\times\id}\\
                &\cP_G(k)\times\A^k\ar[r]_-{\theta_k} & \A
            }
        \end{align*}
        equals
        \begin{align*}
            \xymatrix{
                & \cP_G(k)\times \left(\cP_G(1)\times\A\right)^k \ar[dr]^-{\id\times\theta_1^k} \ar@{=>}[d]^-{\id\times \prod \phi}\\
                \cP_G(k)\times\A^k\ar[ur]^-{\id\times\iota^k} \ar[rr]_-{\id} \ar[d]_-{\id\times\iota^k} \ar[dr]_-{\iota^k} && \cP_G(k)\times\A^k\ar[r]^-{\theta_k} & \A\\
                \cP_G(k)\times\left(\cP_G(1)\times\A\right)^k \ar[r]_-{\text{shuffle}}& \cP_G(k)\times \cP_G(1)^k\times\A^k \ar[ur]_-{\gamma\times\id}
            }
        \end{align*}

        Also:
        \begin{align*}
            \xymatrix{
                \cP_G(k)\times\A^k\ar[r]^-{\iota} \ar[dr]_-{\id} & \cP_G(1)\times\cP_G(k)\times \A^k \ar[r]^-{\id\times\theta_k} \ar[d]_-{\gamma\times\id} & \cP_G(1)\times\A \ar[d]^-{\theta_1} \ar@{=>}[dl]^-{\phi}\\
                &\cP_G(k)\times\A^k \ar[r]_-{\theta_k} & \A
            }
        \end{align*}
        equals
        \begin{align*}
            \xymatrix{
                \cP_G(k)\times \A^k \ar[r]^-{\iota} \ar[d]_-{\theta_k} & \cP_G(1)\times\cP_G(k)\times \A^k \ar[d]^-{\id\times\theta_k}\\
                \A \ar[r]^-{\iota} \ar[dr]_-{\id} & \cP_G(1)\times\A \ar[d]^-{\theta_1} \ar@{=>}[dl]+<30pt>^>>>>>>>{\phi_1}\\
                & \A
            }
        \end{align*}
        The triangles above commute due to the operadic unit axiom.
    \end{itemize}
\end{defn}

\begin{rem}
    In \cite{CG14} Corner and Gurski included unital natural isomorphisms $\phi_1$, while \cite{GMMO20} required strict unitality. The main result in nonequivariant case of \cite{LMRZ26} alleviated the result from strict unitality.
\end{rem}

\begin{defn}[$\cP_G$-pseudomorphism]\label{defn_P_G_pseudomorph}
    A $\cP_G$-pseudomorphism $(\FF,\zeta):(\A,\theta,\phi)\rightarrow(\B,\vartheta,\varphi)$ consists of functor $\FF:\A\rightarrow\B$, such that for each $j$, the following diagram commutes up to an invertible natural transformation $\zeta_j$:
    \begin{align*}
        \xymatrix{
            \cP_G(j)\times \A^j\ar[r]^-{\id\times\FF^j} \ar[d]_-{\theta_j} & \cP_G(j)\times \B^j \ar[d]^-{\vartheta_j} \ar@{=>}[dl]^-{\zeta_j}\\
            \A \ar[r]_-{\FF} & \B
        }
    \end{align*}
    
    Moreover, they satisfy the following properties:
    \begin{itemize}
        \item[(i)]\textbf{(Equivariance of $\zeta$)} For each $j$,$\zeta_j=\zeta_j\circ (\sigma\times\sigma^{-1})$ for all $\sigma\in \Sigma_j$. 
        \item[(ii)] \textbf{(Preserves Identity)} The following two pasting diagrams are equal:
        \begin{align*}
            \xymatrix{
                \A \ar[r]^-{\FF} \ar[d]_-{\iota} & \B \ar[d]_-{\iota} \ar@(r,u)[ddr]^-{\id}\\
                \cP_G(1)\times\A \ar[r]^-{\id\times\FF} \ar[dr]_-{\theta_1} & \cP_G(1)\times\B \ar[dr]_-{\vartheta_1} \ar@{=>}[r]-<10pt>^-{\varphi_1} & \\
                & \A \ar[r]_-{\FF} \ar@{=>}[u]_-{\zeta_1^{-1}} & \B
            }\quad = \quad \xymatrix{
                \A \ar[r]^-{\FF} \ar[d]_-{\iota} \ar[dr]^-{\id} & \B \ar[dr]^-{\id}\\
                \cP_G(1)\times \A \ar@{=>}[ur]-<25pt>^-{\phi_1} \ar[r]_-{\theta_1} & \A \ar[r]_-{\FF} & \B
            }
        \end{align*}
        \item[(iii)] \textbf{(Preserves Composition)} The following two pasting diagrams are equal:
        \begin{align*}
            \xymatrix{
                \cP_G(k)\times \prod_r\left(\cP_G(j_r)\times \A^{j_r}\right) \ar[rr]^-{\id\times\prod(\id\times \FF^{j_r})} \ar[dr]_-{\id\times \prod\theta_{j_r}} \ar[dd]_-{\gamma} && \cP_G(k)\times \prod_r\left(\cP_G(j_r)\times \B^{j_r}\right) \ar[dr]^-{\id\times \prod\vartheta_{j_r}} \ar@{=>}[dl]^-{\id\times \prod \zeta_{j_r}}\\
                & \cP_G(k)\times \A^k \ar[rr]_-{\id\times \FF^k} \ar@{=>}[dl]^-{\phi} \ar[dd]_-{\theta_k} && \cP_G(k)\times \B^k \ar[dd]^-{\vartheta_k} \ar@{=>}[ddll]^-{\zeta_k}\\
                \cP_G(j)\times \A^j \ar[dr]_-{\theta_j}\\
                & \A \ar[rr]_-{\FF} && \B
            }
        \end{align*}
        equals
        \begin{align*}
            \xymatrix{
                \cP_G(k)\times \prod_r\left(\cP_G(j_r)\times\A^{j_r}\right) \ar[rr]^-{\id\times \prod(\id\times \FF^{j_r})} \ar[dd]_-{\gamma} && \cP_G(k)\times \prod\left(\cP_G(j_r)\times \B^{j_r}\right) \ar[dd]_-{\gamma} \ar[dr]^-{\id\times \prod\theta_{j_r}}\\
                &&& \cP_G(k)\times \B^k \ar[dd]^-{\vartheta_k} \ar@{=>}[dl]^-{\varphi}\\
                \cP_G(j)\times\A^j \ar[rr]^-{\id\times \FF^j} \ar[dr]_-{\theta_j} && \cP_G(j)\times \B^j \ar[dr]^-{\vartheta_j} \ar@{=>}[dl]^-{\zeta_j}\\
                & \A \ar[rr]_-{\FF} && \B
            }
        \end{align*}
    \end{itemize}
\end{defn}

\vspace{1pt}

\begin{defn}[$\cP_G$-2-cells]\label{defn_P_G_2_cell}
    Let $(\FF,\zeta),(\GG,\xi):(\A,\theta,\phi)\rightarrow (\B,\vartheta,\varphi)$ be two $\cP_G$-pseudomorphisms, a $\cP_G$-2-cell $\lambda:(\FF,\zeta)\Rightarrow(\GG,\xi)$ is a natural transformation $\lambda:\FF\rightarrow\GG$, such that the following two pasting diagrams are equal:
    \begin{align*}
        \xymatrix{
            \cP_G(j)\times \A^j \ar@(ur,ul)[r]^-{\id\times \FF^j} \ar@{=>}[r]^-{\id\times \lambda^j} \ar@(dr,dl)[r]^-{\id\times\GG^j} \ar[dd]_-{\theta_j} & \cP_G(j)\times \B^j \ar[dd]^-{\vartheta_j} \ar@{=>}[ddl]^-{\xi_j}\\
            \\
            \A \ar[r]_-{\GG} & \B
        }\quad = \quad \xymatrix{
            \cP_G(j)\times \A^j \ar[dd]_-{\theta_j} \ar[r]^-{\id\times \FF^j} & \cP_G(j)\times\B^j \ar[dd]^-{\vartheta_j} \ar@{=>}[ddl]_-{\zeta_j}\\
            \\
            \A \ar@(ur,ul)[r]_-{\FF} \ar@{=>}[r]_-{\lambda} \ar@(dr,dl)[r]_-{\GG} & \B
        }
    \end{align*}
\end{defn}

\begin{rem}\label{rem_gen_G_op}
    The above definitions can be shared across all $G$-operad in $\Cat_G$, it's just the main focus here being $\cP_G$.
\end{rem}

\vspace{1pt}

\section{Category $\F_G$ and its Pseudoalgebras}\label{sec_F_Gm}

Within the equivariant infinite loop space machine, the purpose for $\F_G$- pseudoalgebras (after strictifying to algebras) is to construct $\F_G$-spaces, which are collections of $G$-spaces indexed over finite based $G$-sets (and morphisms between them encoding extra structure). The $\F_G$-spaces modify to genuine $G$-spectra, which the psudoalgebras are the most relevant categorical sources one can study regarding it. For details of $\F_G$-spaces, cf. \cite{S89} Section 1 (which is named after $\Gamma_G$ instead of $\F_G$).

The following definitions are equivariant generalizations of $\F$-pseudofunctors in \cite{LMRZ26}, and we'll utilize similar notations.

\vspace{1pt}

\begin{defn}
    $\F_G$ denotes the \emph{category of based finite $G$-sets}, where the objects are $\mathbf{n}^\alpha$, set wise being the finite set $\{0,1,...,n\}$ with $0$ fixed as the basepoint, and $\alpha:G\rightarrow\Sigma_n$ is a based $G$-action on the set.

    The homset $\F_G(\mathbf{n}^\alpha,\mathbf{m}^\beta)$ consists of all based set maps, which the homset also carries a $G$-action via conjugation. From now on, we'll view $\F_G$ as a 2-category, where each hom-category $\F_G(\mathbf{n}^\alpha,\mathbf{m}^\beta)$ is discrete with only identity 2-cells.

    Furthermore, we define the following:
    \begin{itemize}
        \item $\Pi_G\subset \F_G$ denotes the subcategory where each morphism $\phi:\mathbf{n}^\alpha\rightarrow\mathbf{m}^\beta$ satisfies $|\phi^{-1}(j)| = 0,1$ for all $j\neq 0$.
        \item $\Lambda_G,\Lambda_G^\perp,\Sigma_G\subset\Pi_G$ denotes the subcategory where each morphism is injective, surjective, and bijective respectively (in particular, $\Lambda_G\cap \Lambda_G^\perp = \Sigma_G$).
        \item $\mathscr{E}_G\subset \F_G$ denotes the subcategory where each morphism satisfies $\phi^{-1}(0)=\{0\}$.
    \end{itemize}
    We also denote the map $\phi_{\mathbf{n}^\alpha}:\mathbf{n}^\alpha\rightarrow\mathbf{1}$ (with $j\mapsto 1$ for all $j\neq 0$) in $\mathscr{E}$ as the \emph{canonical product}.
\end{defn}

\vspace{1pt}

\subsection{2-Category $\F_G$-$\PsAlg$}

\begin{defn}[$\F_G$-pseudoalgebra]
    An $\F_G$-pseudoalgebra is a pseudofunctor $\B:\F_G\rightarrow \Cat_G$ (cf. Definition \ref{defn_pseudofunctor}), such that restricting to the subcategory $\Pi_G\subset\F_G$, $\B$ is a strict functor. Moreover, we require it to be reduced, namely $\B(0)=*$ (the terminal $G$-category), and each hom-functor 
    \begin{align*}
        \B_{\mathbf{n}^\alpha,\mathbf{m}^\beta}:\F_G(\mathbf{n}^\alpha,\mathbf{m}^\beta)\rightarrow \Cat_G(\B(\mathbf{n}^\alpha),\B(\mathbf{m}^\beta))
    \end{align*}
    is a $G$-functor.
\end{defn}

\vspace{1pt}

Let $\delta_i:\mathbf{n}^\alpha\rightarrow\mathbf{1}$ be the $i$th projection defined by $i\mapsto 1$, and any $i\neq j\mapsto 0$. Since it belongs to $\Pi_G$, it induces a functor $\B\delta_i:\B(\mathbf{n}^\alpha)\rightarrow \B(\mathbf{1})$. Hence, by universality of product, the collection $\B\delta_i$ ($1\leq i\leq n$) induces a functor $\delta:\B(\mathbf{n}^\alpha)\rightarrow \B(\mathbf{1})^{\mathbf{n}^\alpha}$ (where $\B(\mathbf{1})^{\mathbf{n}}\cong \B(\mathbf{1})^{\mathbf{n}^{\alpha}}$ nonequivariantly).

\begin{defn}[$\F_{G,\vs}$ and $\F_{G,\RR}$-pseudoalgebra]
    An $\F_G$-pseudoalgebra $\B$ is \emph{very special} (or an $\F_{G,\vs}$-pseudoalgebra), if $\delta:\B(\mathbf{n}^\alpha)\rightarrow \B(\mathbf{1})^{\mathbf{n}^\alpha}$ is an isomorphism for all $\mathbf{n}^\alpha\in\F_G$, and $\B$ is \emph{strictly special} (or an $\F_{G,\RR}$-pseudoalgebra), if all $\delta$ is in fact identity.
\end{defn}

\vspace{1pt}

\begin{defn}[$\F_G$-1-cell]
    Given $\A,\B$ two $\F_G$-pseudoalgebras, a morphism (or $\F_G$-1-cell) between them is a pseudotransform $(\FF,\zeta):\A\rightarrow \B$ (cf. Definition \ref{defn_pseudotransform}), whose restriction to $\Pi_G$ is an actual natural transformation. Moreover, it respects $G$-equivariance for every $g\in G$ and morphism $f\in\Mor(\F_G)$:
    \begin{align*}
        \zeta(g\circ f\circ g^{-1}) = g\circ \zeta(f)\circ g^{-1}
    \end{align*}
\end{defn}

\vspace{1pt}

\begin{defn}[$\F_G$-2-cell]
    Given $\tilde{A},\tilde{B}$ two $\F_G$-pseudoalgebras, and two $\F_G$-1-cells $(\FF,\zeta),(\GG,\xi):\tilde{A}\rightarrow\tilde{\B}$, an $\F_G$-2-cell is a modification $\lambda:\FF\Rightarrow\GG$ (cf. Definition \ref{defn_modification}), such that for any morphism $f:\mathbf{n}^\alpha\rightarrow\mathbf{m}^\beta$ in $\F_G$, the following two pasting diagrams are equal:
    \begin{align*}
        \xymatrix{
            \A(\mathbf{n}^\alpha) \ar[dd]_-{\A f} \ar@(ur,ul)[rr]^-{\FF_{\mathbf{n}^\alpha}} \ar@{=>}[rr]^-{\lambda} \ar@(dr,dl)[rr]^-{\GG_{\mathbf{n}^\alpha}} && \B(\mathbf{n}^\alpha) \ar[dd]^-{\B f} \ar@{=>}[ddll]^-{\xi f}\\
            \\
            \A(\mathbf{m}^\beta) \ar[rr]_-{\GG_{\mathbf{m}^\beta}} && \B(\mathbf{m}^\beta)
        }\quad = \quad \xymatrix{
            \A(\mathbf{n}^\alpha) \ar[dd]_-{\B f} \ar[rr]^-{\FF_{\mathbf{n}^\alpha}} && \B(\mathbf{n}^\alpha) \ar[dd]^-{\B f} \ar@{=>}[ddll]_-{\zeta f}\\
            \\
            \A(\mathbf{m}^\beta) \ar@(ur,ul)[rr]_-{\FF_{\mathbf{m}^\beta}} \ar@{=>}[rr]_-{\lambda} \ar@(dr,dl)[rr]_-{\GG_{\mathbf{m}^\beta}} && \B(\mathbf{m}^\beta)
        }
    \end{align*}
\end{defn}

\section{Obstacles in Equivariant Case}\label{sec_results}
The eventual goal is constructing an equivalence between $\F_{G,\RR}$-pseudoalgebra and $\cP_G$-pseudoalgebra (although this is still an open question). \cite{LMRZ26} provides a general construction for the non-equivariant case (or $G=\mathbf{1}$), here we discuss some potential steps and problems when generalizing the procedure to the nonequivariant case.

\subsection{Functor $\cP_G\text{-}\PsAlg\rightarrow\F_{G,\RR}\text{-}\PsAlg$}
Given a $\cP$-pseudoalgebra $\A$, define object map $\B:\F\rightarrow\Cat$ by $\B(\mathbf{n}):=\A^n$. For extending to a pseudofunctor (and restricting to functor on $\Pi$), the central concepts include: 
\begin{itemize}
    \item[(i)] Substituting the vanishing coordinates of maps $\psi:\mathbf{n}\rightarrow\mathbf{m}$ in $\Pi$ using the specified basepoint $\theta_0:*\cong \cP(0)\rightarrow\A$. In particular, since each $j\in\mathbf{m}$ has preimage being empty or a singleton, define $\psi:\A^n\rightarrow\A^m$ by 
    \begin{align*}
        \psi(a_1,...,a_n):= (a_{\psi^{-1}(1)},...,a_{\psi^{-1}(m)}),
    \end{align*}
    with $a_{\psi^{-1}(j)}:=*$ (the basepoint) if $\psi^{-1}(j)=\emptyset$. This defines $\B$ to be a functor with source $\Pi$.
    \item[(ii)] Treating each canonical $n$-fold product $\phi_n:\mathbf{n}\rightarrow\mathbf{1}$ in $\mathscr{E}$ (with $\phi_n(j)=1$ for all $j\neq 0$) using the $n$-fold product $e\times\A\hookrightarrow \cP(n)\times \A\xrightarrow{\gamma}\A$.
    \item[(iii)] Finally, show the decomposition laws, that maps in $\mathscr{E}$ are compositions of permutation and wedge of canonical products (when fixing a specific order, it's unique), while maps in $\F$ decompose into a projection and a map in $\mathscr{E}$. Using the laws to construct functors $\A^n\rightarrow\A^m$ using each map $\psi$.
\end{itemize}

The coherent natural isomorphisms for $\cP$-pseudoalgebra provide the pseudofunctoriality modulo extensive diagram chasing.

\vspace{1pt}

Step (i) generalizes  to the equivariant case, as each $\cP_G$-pseudoalgebra $\A$ has a natural basepoint specified by $\theta_0:*\cong \cP_G(0)\rightarrow\A$, each $\A^{\mathbf{n}^\alpha}\cong \A^n$ nonequivariantly, and the maps in $\Pi_G$ include all nonequivariant ones, hence the same construction carries over. 

For (ii), each canonical product $\phi_{\mathbf{n}^\alpha}:\mathbf{n}^\alpha\rightarrow\mathbf{1}$ has a choice given by $\alpha\times \A^n\hookrightarrow \cP_G(n)\times \A^n\xrightarrow{\theta_n}\A$ by treating $\alpha$ as an object in $\cP_G(n)=\Cat(\mathscr{E}G,\mathscr{E}\Sigma_n)\cong \Sets(G,\Sigma_n)$.

For (iii), the main problem in equivariant case is: Not every map in $\mathscr{E}_G$ can be expressed as a wedge of canonical products, as the preimage of an element may not necessarily be wedges of $G$-orbits. So, potentially some choices of maps need to be made and more combinatorial details are awaiting.

\vspace{1pt}

\subsection{Functor $\cP_G\text{-}\PsAlg\leftarrow\F_{G,\RR}\text{-}\PsAlg$}

For the converse, given an $\F_{\RR}$- pseudofunctor $\B$, define the category $\A:=\B(\mathbf{1})$, then each canonical product $\phi_n:\mathbf{n}\rightarrow\mathbf{1}$ (by $j\mapsto 1$ for all $j\neq 0$) generates an $n$-fold product $\phi_n:\A^n\rightarrow\A$. By appealing to more categrial details, this defines $\A$ as a $\cP$-pseudoalgebra.

The main obstacle in the equivariant case is the non-finiteness of $G$-operad $\cP_G$ when $G\neq \mathbf{1}$: In \cite{BBKO21}, they showed for nontrivial $G$, the operad $\cP_G$ cannot be generated by finite elements, in particular for each $n>3$ there exists at least one object $f\in \cP_G(n)$, that can't be generated using operadic maps, permutation actions, and $G$-actions from the objects in $\cP(j)$ ($j<n$). This imposes some subtlety in assigning the $\cP_G$-action to $\A:= \B(\mathbf{1})$ when $\B$ is an $\F_{G,\RR}$-pseudoalgebra, as the $f$ above doesn't have a natural choice of product $\A^n\rightarrow\A$.

A potential solution that requires verification is restricting to an equivalent suboperad of $\cP_G$, and see whether the equivalence/isomorphism holds for the suboperad.

\vspace{1pt}

\subsection{Other Potential Direction}
In contrast to a direct construction of isomorphism 2-functors as above, another REU participant Gregory Young suggested constructing another \emph{category of operators} that simultaneously contains information about the operad $\cP_G$ and the category $\F_G$.

\vspace{1pt}

\section*{Acknowledgments}  

I'd like to thank my mentor, Jon Peter May, for suggesting this interesting topic to investigate, his wisdom and patience in explaining the general picture of equivariant infinite loop space theory to me, and for the opportunity to participate in this REU, which made this work possible. I'd also like to thank my graduate mentor, Ajay Srinivasan, for having multiple detailed discussions on the categorical properties and the axioms above, which made it possible for me to solidify the underlying categorical groundwork. Finally, I'd like to thank another REU participant, Gregory Young, for suggesting a potential approach to showing an equivalence and some perspectives on studying these abstract objects. 

\begin{thebibliography}{9}

\bibitem{BBKO21}%small change
K. Bangs, S. Binegar, Y. Kim, K. Ormsby, A. M. Osorno, D. Tamas-Parris and L. Xu.
Biased Permutative Equivariant Categories.
Homology, Homotopy and Applications, vol. 23.
2021.
https://arxiv.org/abs/1907.00933

%\bibitem{BP20}
%P. Bonventre, L. A. Pereira.
%Genuine Equivariant Operads.
%Available as ArXiv:0717.02226, 2020.

\bibitem{CG14}
A. S. Corner, N. Gurski.
Operads with general groups of equivariance, and some 2-categorical aspects of operads in $\Cat$.
Available as arXiv:1312.5910v2, 2014.

\bibitem{GM17}
B. J. Guillou and J. P. May.
Equivariant iterated loop space theory and permutative $G$-categories.
\emph{Algebr. Geom. Topol.}, 17(6):3259-3339, 2017.

\bibitem{GMMO18}
B. J. Guillou, J. P. May, M. Merling, A. M. Osorno.
A symmetric monoidal and equivariant segal infinite loop space machine.
Available as arXiv:1711.09183, 2018.

\bibitem{GMMO20}
B. J. Guillou, J. P. May, M. Merling, A. M. Osorno.
Symmetric monoidal $G$-categories and their strictification.
\emph{Q. J. Math.}, 71(1):207-246, 2020.

\bibitem{GMMO21}
B. J. Guillou, J. P. May, M. Merling, A. M. Osorno.
Multiplicative equivariant $K$-Theory and the Barratt-Priddy-Quillen theorem.
Available as arXiv:2102.13246, 2021.

\bibitem{JY21}
N. Johnson, D. Yau.
\emph{2-dimensional categories}.
Oxford University Press, Oxford, 2021.

%\bibitem{KMZ26}
%H. J. Kong, J. P. May, F. Zou.
%Orbital Presheaves in Equivariant Infinite Loop Space Theory.
%Available as ArXiv:2607.12141, 2026.

\bibitem{LMRZ26}
J. Liu, J. P. May, K. I. Roke, H. Zhang and K. Zhou.
What are symmetric monoidal categories?
Available as arXiv:2607.13912, 2026.

\bibitem{M72}
J. P. May. $E_\infty$ spaces, group completions, and permutative categories. In \emph{New Development in Topology (Proc. Sympos. Algebraic Topology, Oxford, 1972)}, pages 61-93. London Math. Soc. Lecture Note Ser., No. 11. Cambridge Univ. Press, London, 1974.

\bibitem{M97}
J. P. May. Operads, algebras and modules. In \emph{Operads: Proceedings of Renaissance Conferences (Hartford, CT/Luminy, 1995)}, volume 202 of \emph{Contemp. Math.}, pages 15-31. Amer. Math. Soc., Providence, RI, 1997. 

\bibitem{M98}
S. Mac Lane.
\emph{Categories for the working mathematician}, volume 5 of \emph{Graduate Texts in Mathematics}. Springer-Verlag, New York, 2nd edition, 1998.

\bibitem{MMO25}
J. P. May, M. Merling, and A. M. Osorno.
Equivariant infinite loop space theory, the space level theory.
Available as arXiv:1704.03413, 2025.

\bibitem{R16}
E. Riehl.
\emph{Category theory in context}. Courier Dover Publications, 2016.

\bibitem{S89}
K. Shimakawa. 
Infinite loop $G$-spaces associated to monoidal $G$-graded categories.
Publ. Res. Inst. Math. Sci 25, no. 2, pp. 239-262, 1989.

\bibitem{Y24}
D. Yau. Multifunctorial equivariant algebraic $K$-theory. Available as arXiv:2404.02794, 2024.

\end{thebibliography}

\end{document}